\documentclass[12pt, a4paper]{article}

% ===== Encoding and fonts =====
\usepackage[utf8]{inputenc}
\usepackage[T1]{fontenc}
\usepackage{lmodern}
\usepackage{microtype}

% ===== Mathematics =====
\usepackage{amsmath}
\usepackage{amssymb}
\usepackage{amsthm}
\usepackage{mathtools}
\usepackage{thmtools}

% ===== Geometry =====
\usepackage[
  top=2.5cm,
  bottom=2.5cm,
  left=3cm,
  right=3cm
]{geometry}

% ===== References and links =====
\usepackage[
  backend=biber,
  style=alphabetic,
  sorting=nyt,
  maxnames=4
]{biblatex}
\addbibresource{references.bib}

\usepackage[
  colorlinks=true,
  linkcolor=blue!60!black,
  citecolor=green!50!black,
  urlcolor=blue!60!black
]{hyperref}
\usepackage{cleveref}

% ===== Theorem environments =====
\declaretheoremstyle[
  headfont=\bfseries,
  notefont=\normalfont,
  bodyfont=\itshape,
  headpunct={.},
  postheadspace=0.5em
]{thmstyle}

\declaretheoremstyle[
  headfont=\bfseries,
  notefont=\normalfont,
  bodyfont=\normalfont,
  headpunct={.},
  postheadspace=0.5em
]{defstyle}

\declaretheorem[style=thmstyle, name=Theorem,     numberwithin=section]{theorem}
\declaretheorem[style=thmstyle, name=Lemma,       sibling=theorem]{lemma}
\declaretheorem[style=thmstyle, name=Proposition, sibling=theorem]{proposition}
\declaretheorem[style=thmstyle, name=Corollary,   sibling=theorem]{corollary}
\declaretheorem[style=defstyle, name=Definition,  sibling=theorem]{definition}
\declaretheorem[style=defstyle, name=Remark,      sibling=theorem]{remark}
\declaretheorem[style=defstyle, name=Example,     sibling=theorem]{example}
\declaretheorem[style=defstyle, name=Conjecture,  sibling=theorem]{conjecture}
\declaretheorem[style=defstyle, name=Heuristic,   sibling=theorem]{heuristic}

% ===== Colors and notes =====
\usepackage{xcolor}
\usepackage{booktabs}
\newcommand{\gap}[1]{%
  \textcolor{red!70!black}{\textbf{[Bridge requiring proof:} #1\textbf{]}}%
}
\newcommand{\todo}[1]{%
  \textcolor{orange!80!black}{\textbf{[TODO:} #1\textbf{]}}%
}
\newcommand{\openq}[1]{%
  \textcolor{blue!70!black}{\textbf{[Open:} #1\textbf{]}}%
}

% ===== Operators and notation =====
\newcommand{\C}{\mathbb{C}}
\newcommand{\R}{\mathbb{R}}
\newcommand{\N}{\mathbb{N}}
\newcommand{\Z}{\mathbb{Z}}
\renewcommand{\H}{\mathbb{H}}
\renewcommand{\Re}{\operatorname{Re}}
\renewcommand{\Im}{\operatorname{Im}}
\newcommand{\cl}{\mathcal{L}}
\newcommand{\hs}{\mathcal{H}}

% ===== Title metadata =====
\title{%
  \textbf{Riemann Hypothesis via Spectral Coherence}\\[0.5em]
  \large A Proposed Framework Linking Reflection Symmetry\\
  to the Critical Line
}

\author{%
  \textit{[Author name(s) to be added]}\\[0.3em]
  \small\textit{[Institution / affiliation]}\\
  \small\textit{[Contact address]}
}

\date{%
  Version 0.1.0-draft\\
  \today\\[0.5em]
  \small\textit{Manuscript under active development and review.}
}

\begin{document}

\maketitle

\begin{abstract}
This manuscript develops a proposed proof framework for the Riemann Hypothesis through
a spectral-coherence interpretation of the functional symmetry of the Riemann zeta
function $\zeta(s)$. The central object of study is the completed zeta function
$\xi(s) = \frac{1}{2}s(s-1)\pi^{-s/2}\Gamma(s/2)\zeta(s)$, which satisfies the
functional equation $\xi(s) = \xi(1-s)$ and is symmetric under reflection about
the critical line $\Re(s) = \tfrac{1}{2}$.

The manuscript proposes a spectral-coherence structure associated with $\xi$ in
which the nontrivial zeros of $\zeta$ appear as \emph{balance points} — states at
which a prescribed cancellation condition is satisfied. The framework attempts to
isolate the balance mechanism that would force all nontrivial zeros onto the
critical line $\Re(s) = \tfrac{1}{2}$, by showing that the relevant balance
condition is structurally incompatible with $\Re(s) \ne \tfrac{1}{2}$.

Specifically, the manuscript aims to:
\begin{enumerate}
  \item Define a spectral operator $H$ whose spectral properties encode the
        symmetry of $\xi$;
  \item Establish that nontrivial zeros of $\zeta$ correspond exactly to balance
        points of $H$;
  \item Prove that all balance points of $H$ lie on $\Re(s) = \tfrac{1}{2}$.
\end{enumerate}
Steps~(2) and~(3) constitute the primary open problems in the current version of
the manuscript. The framework is presented for mathematical review, with all
unresolved steps explicitly identified.

\medskip
\noindent\textbf{Keywords:} Riemann Hypothesis; zeta function; spectral theory;
reflection symmetry; critical line; coherence; operator theory.

\end{abstract}

\vspace{1em}
\noindent\textbf{Disclaimer.}
This manuscript presents a proposed proof framework and is offered for mathematical
review, critique, and formalization. It does not claim to be a complete proof of
the Riemann Hypothesis. All steps that have not been rigorously established are
explicitly marked.

\tableofcontents
\newpage

% ===== Main sections =====
\section{Introduction}
The Riemann Hypothesis (RH) states that every nontrivial zero of $\zeta(s)$ has real part $\tfrac12$.
This manuscript attempts a direct reduction built on one explicit function:
\[
\Phi(s):=\Gamma(s)(1-s)\zeta(s),\qquad s\in\C.
\]

The proof spine has three steps.
First, in the critical strip $0<\RePart(s)<1$, show
\[
\Phi(s)=0\iff\zeta(s)=0.
\]
This is an elementary correspondence because $\Gamma(s)$ and $1-s$ are nonzero in that domain.

Second, isolate the only substantial theorem target:
\[
\Phi(s)=0\Longrightarrow \RePart(s)=\tfrac12
\quad (0<\RePart(s)<1).
\]
This implication is the central open bridge in the current draft.

Third, note the conditional reduction: if the implication above is proved, then RH follows immediately from the strip correspondence.

The manuscript is intentionally audit-first.
It distinguishes established setup from open burden, and it does not present heuristic or computational support as a substitute for proof.


\section{Definitions and Notation}
This section fixes baseline notation and separates classical definitions from manuscript-specific provisional objects.

\subsection*{Classical definitions (standard)}

\begin{definition}[Completed zeta function]
\[
\xi(s) := \tfrac12 s(s-1)\pi^{-s/2}\Gamma(s/2)\zeta(s).
\]
We use the standard functional-equation symmetry $\xi(s)=\xi(1-s)$.
\end{definition}

\begin{definition}[Critical strip and critical line]
The critical strip is
\[
\{s\in\C:0<\RePart(s)<1\},
\]
and the critical line is
\[
\{s\in\C:\RePart(s)=\tfrac12\}.
\]
\end{definition}

\begin{definition}[Reflection involution]
Let $\sigma:\C\to\C$ be $\sigma(s)=1-s$. This map is an involution, and its fixed-point set is the critical line.
\end{definition}

\subsection*{Manuscript-specific provisional definitions}

\begin{definition}[Proposed coherence object $H$ \textnormal{(provisional definition; placeholder object)}]
\placeholder{A full operator-level construction is still required: ambient space, domain, action, and well-posedness properties.}
The symbol $H$ denotes a proposed spectral object used by the manuscript as a formalization target for the coherence program.
At the current stage, $H$ is reserved notation, not yet a fully constructed theorem-level object.
\end{definition}

\begin{definition}[Admissible point and admissible class \textnormal{(provisional definition; formalization target)}]
\placeholder{A full well-posedness theorem for the regularized trace below is still required.}
Fix the proposed coherence object $H$.
An \emph{admissible point} is a complex number $s\in\C$ such that both regularized resolvents
\[
(H-sI)^{-1}_{\mathrm{reg}},
\qquad
(H-(1-s)I)^{-1}_{\mathrm{reg}}
\]
exist in the chosen operator model and have finite regularized traces.
The \emph{admissible class} is
\[
\mathcal{A}:=\left\{s\in\C:
\Tr_{\mathrm{reg}}(H-sI)^{-1}\ \text{and}\ \Tr_{\mathrm{reg}}(H-(1-s)I)^{-1}\ \text{are both defined and finite}
\right\}.
\]
\end{definition}

\begin{definition}[Canonical balance functional \texorpdfstring{$\Phi_H$}{Phi_H} \textnormal{(single main-text candidate; conditional/provisional)}]
\placeholder{The regularization scheme and its theorem-level properties remain open proof obligations.}
For $s\in\mathcal{A}$, define
\[
\Phi_H(s)
:=
\Tr_{\mathrm{reg}}(H-sI)^{-1}
-
\Tr_{\mathrm{reg}}(H-(1-s)I)^{-1}.
\]
Thus $\Phi_H:\mathcal{A}\to\C$ is the \emph{regularized resolvent reflection-defect functional} of $H$.
The manuscript uses the single balance predicate
\[
\text{Bal}_H(s) \iff \Phi_H(s)=0 \qquad (s\in\mathcal{A}).
\]
Interpretation: $\Phi_H(s)=0$ encodes exact reflection balance of the regularized resolvent response at the pair $(s,1-s)$.
This is the canonical main-text candidate for the balance object, and it remains a theorem target in the current manuscript version.
\end{definition}

\begin{remark}[Why this is the canonical candidate in the main flow]
The functional above is used as the single main-text candidate because it directly measures the $s\leftrightarrow 1-s$ defect at operator level, which is the symmetry direction needed by the proof spine.
No alternative balance functional is used in the core manuscript flow.
\end{remark}

\begin{remark}[Status boundary for theorem-level use]
The definitions of $H$, admissibility, and $\Phi_H$ above are provisional.
Open bridges include: construction of the operator model, choice and invariance of regularization, existence/finite-value proofs on $\mathcal{A}$, and compatibility of this $\Phi_H$ with the zero-correspondence target.
No theorem-level claim may rely on these terms without first supplying those formal components.
\end{remark}

\begin{remark}[Traceability requirement]
When the provisional objects are upgraded, each promoted definition should include: (i) explicit hypotheses, (ii) well-posedness statement, and (iii) dependency links to the lemmas/theorems that justify subsequent use.
\end{remark}

\subsection*{Stage-6 provisional transport-defect interface}

\begin{definition}[No-drift defect functional \texorpdfstring{$\Delta_{\mathrm{nd}}$}{Delta\_nd} \textnormal{(provisional; audit/planning object)}]
Work under one synchronized package $\mathcal{H}_{\mathrm{S6}}^{\min}$, and fix:
\begin{enumerate}
  \item a Stage-5 hypothesis tuple $\Theta_{\mathrm{S5}}=(\theta_1,\dots,\theta_m)$;
  \item a transfer class $\mathfrak{C}_{\mathrm{tr}}$ and contradiction-relevant subclass
  $\mathfrak{C}_{\mathrm{tr}}^{\mathrm{ctr}}\subseteq\mathfrak{C}_{\mathrm{tr}}$;
  \item for each $C\in\mathfrak{C}_{\mathrm{tr}}$, a transfer-position set $U(C)\subseteq\mathcal{A}_{\mathrm{strip}}$;
  \item for each clause $\theta_j$ and each $C$, a clause-defect gauge
  $\delta_j^{(C)}:U(C)\to[0,\infty]$.
\end{enumerate}
Define the pointwise tuple-defect gauge
\[
d_{\Theta}^{(C)}(u):=\max_{1\le j\le m}\delta_j^{(C)}(u),
\]
and define the no-drift defect of $C$ by
\[
\Delta_{\mathrm{nd}}(C;\Theta_{\mathrm{S5}},\delta):=
\sup_{u\in U(C)}d_{\Theta}^{(C)}(u)\in[0,\infty].
\]
Thus $\Delta_{\mathrm{nd}}$ is a \emph{transport-instability measure}: a supremum defect over positions $u$ and clauses $\theta_j$ along a transfer object $C$.
Intended semantics (not proved): $\Delta_{\mathrm{nd}}(C)=0$ means exact clause preservation along $C$; $0<\Delta_{\mathrm{nd}}(C)\ll1$ means quantitatively small drift; finite uniform bounds on a transfer subclass mean bounded instability; $\Delta_{\mathrm{nd}}(C)=\infty$ means uncontrolled transport.
\end{definition}

\begin{remark}[Status and dependency boundary]
This definition is provisional and interface-level.
It does not select a finalized operator realization and does not assert that the chosen gauges are canonical.
Its current role is to provide a clean first analytic control target for \cref{lem:stage6-first-concrete-no-drift}, downstream of Stage-5 local uniqueness and upstream of \cref{thm:stage6-neutralizing-transport-uniformity}.
\end{remark}


\section{Symmetry Framework}
This section records only classical consequences of the completed-zeta symmetries.
It is a boundary statement: these facts constrain admissible zero patterns, but do not by themselves prove RH.

Let
\[
\xi(s)=\tfrac12 s(s-1)\pi^{-s/2}\Gamma\!\left(\tfrac{s}{2}\right)\zeta(s).
\]
The standard functional equation is
\[
\xi(s)=\xi(1-s),
\]
and $\xi(\bar s)=\overline{\xi(s)}$ (equivalently, $\xi$ is real on the real axis) \cite{Edwards1974,Titchmarsh1986}.
Define the reflection map $\sigma(s)=1-s$.

\begin{proposition}[Classical symmetry consequences]
\label{prop:symmetry-boundary}
Let $\rho$ be a nontrivial zero of $\zeta$ (equivalently, of $\xi$). Then:
\begin{enumerate}
  \item $\sigma(\rho)=1-\rho$ is also a zero.
  \item $\bar\rho$ is also a zero.
  \item Hence the zero set is invariant under the group generated by $s\mapsto 1-s$ and $s\mapsto \bar s$; in particular, zeros off the critical line and off the real axis occur in quartets
  \[
  \{\rho,\,1-\rho,\,\bar\rho,\,1-\bar\rho\}.
  \]
\end{enumerate}
These are standard consequences of the functional equation and conjugation symmetry \cite{Edwards1974,Titchmarsh1986}.
\end{proposition}

\begin{remark}[Boundary of what symmetry alone yields]
Proposition~\ref{prop:symmetry-boundary} provides pairing and orbit structure only.
It does \emph{not} exclude the possibility of symmetric off-line zeros $\rho$ with $\RePart(\rho)\neq \tfrac12$.
Therefore any argument for RH requires an additional exclusion mechanism beyond functional-equation symmetry.
\end{remark}

\auditnote{This section is intentionally limited to established symmetry facts; no manuscript-specific mechanism is invoked here.}


\section{Spectral Framework}
This section fixes one canonical main-text answer to the question ``What is $\Phi_H(s)$?''
It does not claim this choice is already validated.
Its role is to state the exact object and list the theorem burdens required before RH-reduction claims become theorem-level statements.

\subsection*{Canonical main-text choice}
\placeholder{The operator construction and regularization theorem package are not yet complete.}
The manuscript fixes one functional in the main flow:
\[
\Phi_H(s):=\Tr_{\mathrm{reg}}(H-sI)^{-1}-\Tr_{\mathrm{reg}}(H-(1-s)I)^{-1},
\]
defined on
\[
\mathcal{A}:=
\left\{
s\in\C:
\Tr_{\mathrm{reg}}(H-sI)^{-1},\Tr_{\mathrm{reg}}(H-(1-s)I)^{-1}\ \text{exist and are finite}
\right\}.
\]
The associated predicate is
\[
\text{Bal}_H(s)\iff \Phi_H(s)=0 \qquad (s\in\mathcal{A}).
\]
Interpretation: $\Phi_H(s)$ is the regularized resolvent reflection defect of $H$ at the pair $(s,1-s)$.

\subsection*{Minimal required structure for the $H$-layer}
The manuscript does not need every possible operator-theoretic feature.
It does need the following \emph{minimal interface} to keep later arguments well-typed and auditable.

\paragraph{(M1) Ambient space and typing data.}
A specified complex topological vector space (typically a Hilbert or Banach model) $\mathcal{X}$, together with clear codomain conventions for all derived maps.
Every use of $H$ must state its typing context explicitly.

\paragraph{(M2) Domain and well-defined action.}
A dense or otherwise explicitly justified domain $\mathcal{D}(H)\subseteq\mathcal{X}$ (as required by the chosen model), and a map
\[
H:\mathcal{D}(H)\to\mathcal{X}
\]
with proof that this action is well defined under stated hypotheses.
If closure/closability is needed by subsequent lemmas, that requirement must be stated and proved where used.

\paragraph{(M3) Symmetry-compatibility package.}
A precise compatibility statement linking the operator layer to the reflection involution $\sigma(s)=1-s$ from \texttt{definitions.tex}.
This must be explicit and non-heuristic.
Informal symmetry pictures are admissible only as motivation, not as substitute for this package.

\paragraph{(M4) Operator-regularity assumptions (structural side).}
The exact regularity required for operator statements: continuity/measurability/analytic dependence of operator-valued maps, domain persistence, and closure/graph conditions needed for later lemmas.
Each assumption must be tagged as either proved in-manuscript or imported as an explicit hypothesis.

\paragraph{(M5) Dependency traceability.}
A lemma-level dependency map from operator properties to later claims (especially Targets~A and~B).
No RH-relevant step should rely on an unlisted implicit property of $H$.

\subsection*{Required theorem package for this explicit \texorpdfstring{$\Phi_H$}{Phi_H}}
The definition above is explicit; the open work is to prove it is mathematically sound and suitable for Targets~A/B.

\paragraph{(P1) Non-circular construction.}
The operator model and regularization map $\Tr_{\mathrm{reg}}$ must be defined without using the zero set of $\zeta$.

\paragraph{(P2) Domain of evaluation (admissible class).}
The declared set $\mathcal{A}$ must be proved nonempty and explicitly characterized.
If singular points are removed, the exclusion mechanism must be stated.

\paragraph{(P3) Well-posedness and regularity.}
For $s\in\mathcal{A}$, prove existence and finiteness of both regularized traces and of their difference.
State the continuity/analytic properties of $s\mapsto\Phi_H(s)$ used later.

\paragraph{(P4) Reflection behavior.}
The explicit formula gives
\[
\Phi_H(1-s)=-\Phi_H(s)
\]
whenever both sides are defined.
A theorem-level proof must show that the admissible class is stable under $s\mapsto 1-s$ and that the regularization is compatible with this identity.

\paragraph{(P5) Bridge to zero claims.}
A precise theorem target must explain how
\[
\Phi_H(s)=0
\]
for this explicit resolvent-defect functional maps to nontrivial zero information.
The bridge should separate:
\begin{enumerate}
  \item the forward direction to be proved,
  \item the converse direction to be proved,
  \item side conditions preventing spurious balanced points.
\end{enumerate}

\subsection*{Canonical exclusion target supported by this interface}
Let
\[
\mathcal{A}_{\mathrm{strip}}:=\{s\in\mathcal{A}:0<\RePart(s)<1\},
\qquad
\mathcal{B}:=\{s\in\mathcal{A}_{\mathrm{strip}}:\Phi_H(s)=0\}.
\]
The main-text exclusion burden is the single subset statement
\[
\mathcal{B}\subseteq\{s\in\mathcal{A}_{\mathrm{strip}}:\RePart(s)=\tfrac12\}.
\]
This section supports that statement by listing operator and analytic prerequisites.
It does not claim those prerequisites are proved.

\subsection*{Open support packages (dependency statements, not proved results)}
\paragraph{Operator-side package (OP).}
The following remain open theorem targets: well-defined operator action and closure properties, existence of the regularized resolvent traces, regularization invariance/consistency, and dependency-locked usage of these results across Targets~A/B.

\paragraph{Analytic-control package (AC).}
The following remain open theorem targets: strip-wide bound control for the explicit $\Phi_H$, legitimacy of limit/interchange operations for regularized traces, continuation control from local to strip-wide regimes, and exclusion of hidden off-line branches of $\mathcal{B}$ caused by analytic degeneration.

\paragraph{Status note.}
OP and AC are recorded as open bridge classes for audit.
They specify proof burden only.
They do not provide a proof of Target~A or Target~B.

\subsection*{Why this explicit choice is kept in the main flow}
The conditional reduction step in \texttt{main-results.tex} depends on theorem-level statements about the specific functional above.
Without a rigorous foundation for this explicit $\Phi_H$:
\begin{enumerate}
  \item $\Phi_H(s)=0$ is not yet a theorem-level condition on a proved domain,
  \item the zero--balance correspondence target lacks a justified bridge,
  \item the canonical exclusion target lacks a proved analytic object to control.
\end{enumerate}
Therefore this explicit $\Phi_H$ choice is a prerequisite object for the manuscript's proposed RH reduction route.

\begin{heuristic}
Geometric or cancellation intuition can guide candidate constructions of $H$ and $\Phi_H$.
However, such intuition remains heuristic until converted into explicit definitions, hypotheses, and proved lemmas matching the interface above.
\end{heuristic}


\section{Lemmas}
We collect auxiliary lemmas used in the main argument. Lemmas that are proved in
the manuscript are marked (Proved). Lemmas that are asserted but not yet proved
are marked with \gap{} and should be treated as conjectures.

\subsection*{Lemmas on the Symmetry Framework}

\begin{lemma}[Reflection Preserves Zeros]
  \label{lem:reflection-zeros}
  If $\zeta(s_0) = 0$ for $s_0$ in the critical strip, then $\zeta(1 - s_0) = 0$.
\end{lemma}

\begin{proof}
  This follows immediately from the functional equation $\xi(s) = \xi(1-s)$ and
  the definition of $\xi$. Since $\xi(s_0) = 0$ and $\xi(s_0) = \xi(1-s_0)$,
  we have $\xi(1-s_0) = 0$, hence $\zeta(1-s_0) = 0$ (the extra factors in $\xi$
  have no zeros in the critical strip).
\end{proof}

\begin{lemma}[Symmetry Under Complex Conjugation]
  \label{lem:conj-zeros}
  If $\zeta(s_0) = 0$, then $\zeta(\bar{s}_0) = 0$.
\end{lemma}

\begin{proof}
  Standard: $\zeta(\bar{s}) = \overline{\zeta(s)}$ for $s$ not equal to $1$.
  Hence $\zeta(\bar{s}_0) = \overline{\zeta(s_0)} = \bar{0} = 0$.
\end{proof}

\begin{lemma}[Zeros Come in Quadruples or Lie on Symmetry Axes]
  \label{lem:quadruples}
  Each nontrivial zero $s_0$ belongs to the set $\{s_0, 1-s_0, \bar{s}_0, 1-\bar{s}_0\}$.
  This set has four distinct elements unless $s_0$ lies on the critical line
  $\Re(s_0) = \tfrac{1}{2}$ or on the real axis $\Im(s_0) = 0$.
\end{lemma}

\begin{proof}
  Combine Lemmas~\ref{lem:reflection-zeros} and~\ref{lem:conj-zeros}. The four
  elements coincide in pairs when $s_0 = 1 - s_0$ (i.e., $\Re(s_0) = \tfrac{1}{2}$)
  or when $s_0 = \bar{s}_0$ (i.e., $\Im(s_0) = 0$, which does not occur for
  nontrivial zeros by standard arguments).
\end{proof}

\subsection*{Lemmas on the Spectral Object}

\begin{lemma}[Symmetry of the Operator]
  \label{lem:op-symmetry}
  \gap{To be proved in \S4 once $H$ is fully defined.}
  The spectral coherence operator $H$ satisfies $H \sim H \circ \sigma$, where
  $\sim$ denotes unitary equivalence and $\sigma \colon s \mapsto 1-s$.
\end{lemma}

\begin{lemma}[Balance Criterion is Well-Posed]
  \label{lem:balance-well-posed}
  \gap{To be proved in \S4.}
  The balance criterion (Definition~\ref{def:balance}) is well-defined: for each
  $s$ in the critical strip, the criterion either holds or fails, and the set of
  $s$ satisfying it is closed.
\end{lemma}

\begin{lemma}[Balance Is Symmetric Under Reflection]
  \label{lem:balance-symmetric}
  \gap{To be proved once Lemma~\ref{lem:op-symmetry} and
  Lemma~\ref{lem:balance-well-posed} are established.}
  If $s$ satisfies the balance criterion, then so does $\sigma(s) = 1-s$.
\end{lemma}

\subsection*{Lemmas on Off-Line Behavior}

\begin{lemma}[Off-Line Asymmetry]
  \label{lem:offlineasymmetry}
  \gap{This is the key lemma required for Step~4 of the proof architecture.
  It must be proved that for $s$ with $\Re(s) \ne \tfrac{1}{2}$, the balance
  criterion fails. This lemma is currently open and is the primary research
  objective.}
  For all $s$ in the critical strip with $\Re(s) \ne \tfrac{1}{2}$, $s$ does
  not satisfy the balance criterion.
\end{lemma}


\section{Main Results}
This section separates established inputs from draft theorem targets.

\subsection*{Established inputs (from standard references)}
\begin{proposition}[Functional symmetry]
The completed zeta function satisfies $\xi(s)=\xi(1-s)$.
\end{proposition}

\begin{proposition}[Zero-set symmetries]
Nontrivial zeros are stable under $s\mapsto 1-s$ and complex conjugation.
\end{proposition}

\subsection*{Draft theorem targets (not proved in this manuscript version)}
\begin{drafttheorem}[Target A: zero--balance correspondence]
\label{thm:target-correspondence}
\placeholder{Draft theorem target. Proof not yet provided.}
Assuming a completed definition of $H$ and the balance criterion, nontrivial zeros in the critical strip correspond exactly to balanced points.
\end{drafttheorem}

\begin{drafttheorem}[Target B: uniqueness of balance locus]
\label{thm:target-uniqueness}
\placeholder{Draft theorem target. Proof not yet provided.}
Balanced points satisfy $\RePart(s)=\tfrac12$.
\end{drafttheorem}

\begin{conjecture}[RH reduction under targets A and B]
If Targets A and B are established under their stated assumptions, RH follows.
\end{conjecture}


\section{Proof Architecture}
We use a four-step proof spine.
The purpose is to keep the main argument transparent for mathematical audit.

\subsection*{Main proof spine}
\begin{enumerate}
  \item \textbf{Define balance.}
  Fix $\mathcal{A}$, define $\Phi_H$, and define
  \[
  \mathcal{B}=\{s\in\mathcal{A}_{\mathrm{strip}}:\Phi_H(s)=0\}.
  \]

  \item \textbf{Target A (correspondence).}
  Prove that nontrivial zeros of $\zeta$ correspond exactly to points in $\mathcal{B}$
  (\cref{thm:target-correspondence}).

  \item \textbf{Target B (critical-line exclusion).}
  Prove that all points of $\mathcal{B}$ satisfy $\RePart(s)=\tfrac12$
  (\cref{thm:stage6-master-exclusion-target}).

  \item \textbf{Deduce RH.}
  Combine Steps 2--3 to obtain the conditional RH reduction
  (\cref{prop:rh-reduction-main-spine}).
\end{enumerate}

\subsection*{Status}
Step~1 is definitional.
Steps~2 and~3 are open theorem targets in this manuscript version.
Therefore Step~4 is conditional, and RH is not claimed as proved here.

\subsection*{Where technical machinery now lives}
Detailed transfer frameworks, defect-function constructions, and Stage~6 subroutines are retained as technical supporting material (appendices/notes), not as part of the main proof spine. In particular, the prior long OP/AC cross-cutting catalog is no longer in the main proof architecture; only a short dependency acknowledgment remains in manuscript-level summaries.


\section{The Critical Line Argument}
This section records the manuscript's conditional reduction step in explicit dependency form.
The purpose is to isolate the final logical deduction from the earlier bridges that still require proof.

\subsection*{Conditional reduction statement}
\begin{proposition}[Framework-internal conditional RH reduction]
Assume the manuscript definitions of the state space $H$, admissibility class, and balance criterion are fixed and well-posed under a hypothesis package $\mathcal{A}$.
Assume further that the following three inputs hold under the same $\mathcal{A}$:
\begin{enumerate}
  \item \textbf{Target A (zero--balance correspondence).} Every nontrivial zero in the critical strip corresponds to a balanced point in the admissible class (with converse direction as specified in the correspondence theorem).
  \item \textbf{Target B (critical-line localization of balanced points).} Every admissible balanced point satisfies $\RePart(s)=\tfrac12$.
  \item \textbf{Global off-line exclusion package.} No admissible off-line balanced configuration exists in the full critical strip, including boundary/degenerate regimes required by the framework's parameterization.
\end{enumerate}
Then every nontrivial zero $\rho$ of $\zeta$ satisfies $\RePart(\rho)=\tfrac12$.
\end{proposition}

\subsection*{Separation of burdens: local uniqueness vs global exclusion}
For audit clarity, the manuscript treats two logically distinct obligations.
\begin{enumerate}
  \item \textbf{Local uniqueness-of-balance burden (Target B core).}
  Prove that in a specified regime $\mathcal{R}_{\mathrm{loc}}\subseteq\mathcal{A}$, balance implies $\RePart(s)=\tfrac12$.
  This is a localization statement.
  \item \textbf{Global off-line exclusion burden (Stage 6 core).}
  Prove that no admissible balanced point with $\RePart(s)\neq\tfrac12$ exists anywhere in the full strip region of interest, including non-perturbative, boundary-near, and large-parameter regimes not automatically covered by $\mathcal{R}_{\mathrm{loc}}$.
\end{enumerate}
The second item is not a restatement of the first.
It is an extension/closure problem requiring additional uniform control.

\subsection*{Why symmetry alone is insufficient}
Classical symmetries give transport rules ($s\mapsto1-s$, $s\mapsto\bar s$).
They do not by themselves imply that the only balance-supporting locus is $\RePart(s)=\tfrac12$.
An argument can preserve symmetry and still allow paired off-line candidates.
Hence the framework requires a separate exclusion mechanism beyond symmetry compatibility.

\subsection*{Reduction proof skeleton (conditional)}
The deduction itself is short once the above inputs are available.
Let $\rho$ be a nontrivial zero.
By Target A, $\rho$ maps to a balanced admissible configuration.
By Target B, that balanced configuration must lie on $\RePart(s)=\tfrac12$.
Hence $\RePart(\rho)=\tfrac12$.
The global off-line exclusion package ensures that no admissible off-line cases survive outside the local regime used to prove Target B, so the conclusion applies across the entire critical strip addressed by $\mathcal{A}$.

\subsection*{Where the argument is formally straightforward}
\begin{itemize}
  \item The final implication from ``zero is balanced'' plus ``balanced implies critical line'' to RH-in-strip is a direct composition of statements.
  \item No deep analytic estimate enters at this final step once hypotheses are synchronized and side conditions are already checked.
\end{itemize}

\subsection*{Where the unresolved burden lies}
The mathematical pressure is upstream, not in the final deduction.
The unresolved burden is to prove, under explicit assumptions, the six core open bridges that feed the reduction:
\begin{enumerate}
  \item \textbf{Exact spectral object burden:} finalize the $H$-layer with explicit domain/action/symmetry data and proof-level well-posedness.
  \item \textbf{Zero--balance correspondence burden (Target A):} establish a rigorous correspondence map with precise domain/codomain control, including both directions and exclusion of spurious balanced states.
  \item \textbf{Uniqueness-of-axis burden (Target B):} prove a theorem-level localization statement forcing every admissible balanced point onto $\RePart(s)=\tfrac12$.
  \item \textbf{Off-critical-line exclusion burden:} upgrade local or regime-restricted arguments to a full-strip exclusion theorem compatible with the assumptions used in Targets A and B.
  \item \textbf{Operator-theoretic burden:} prove each operator property actually used in Targets A/B and exclusion (no hidden assumptions).
  \item \textbf{Analytic-control burden:} justify every growth/convergence/interchange step needed to pass from local constructions to strip-wide statements.
\end{enumerate}
Until these bridges are proved, the reduction remains conditional.

\subsection*{Insufficient argument patterns (explicitly excluded)}
The manuscript does not treat the following as adequate for Target~B or global exclusion:
\begin{enumerate}
  \item \textbf{Pure symmetry arguments} that establish only mirrored pairing of candidate points.
  \item \textbf{Local-only coercivity} without a theorem transferring estimates to the full admissible strip.
  \item \textbf{Regime-fragmented proofs} whose hypotheses differ across subregions without compatibility/patching lemmas.
  \item \textbf{Numerical-only exclusion} not converted into theorem-level analytic statements.
\end{enumerate}

\subsection*{Assumption-compatibility requirements}
For the reduction to be valid as stated, the assumptions used in Target A, Target B, and global off-line exclusion must be mutually compatible.
In particular, the manuscript must avoid: (i) proving Target A on one admissible class and Target B on a narrower incompatible class without a transfer theorem, and (ii) invoking a global exclusion principle whose hypotheses are not inherited by the correspondence setup.

\begin{remark}[Audit-facing interpretation]
This section does \emph{not} claim that RH is proved.
It states that the framework yields RH \emph{if} Targets A and B and the global off-line exclusion package are established under a common explicit hypothesis set.
The hard work is proving those prior bridges.
\end{remark}

\begin{remark}[Common failure modes in RH-type reductions]
Two recurring failure modes are especially relevant here.
First, a correspondence theorem may hold only with hidden regularity assumptions, leaving uncovered zero candidates.
Second, a localization argument may be valid only in a perturbative or local regime and fail to imply global off-line exclusion.
The present framework therefore keeps Target A, Target B, and global exclusion as separate proof obligations.
\end{remark}

\openbridge{Current manuscript status: the conditional deduction is structurally clear, but the six core bridges (exact spectral object, zero--balance correspondence, uniqueness of the balance axis, off-critical-line exclusion, operator-theoretic closure, analytic-control closure) remain open proof obligations requiring full theorem-level proofs under synchronized assumptions.}


\section{Discussion and Related Work}
This discussion section is interpretive.
It is not part of the formal proof itself.

\paragraph{Purpose of the current revision.}
The manuscript now presents one short proof spine in the main text:
(1) define one canonical balance predicate $\Phi_H(s)=0$,
(2) state correspondence,
(3) state one canonical exclusion target,
(4) derive RH conditionally.
This is a presentation and scoping decision for audit clarity.
It does not reduce the technical burden of the open steps.

\paragraph{What the framework attempts to show.}
The argument attempts to show that if the canonical balance predicate is put in exact correspondence with nontrivial zeros, and if the canonical exclusion statement
\[
\mathcal{B}\subseteq\{s\in\mathcal{A}_{\mathrm{strip}}:\RePart(s)=\tfrac12\}
\]
is proved, then RH follows.
The framework suggests a spectral-coherence route to those two targets, but those targets remain open in the current manuscript version.

\paragraph{Why mathematical branching was reduced in the main flow.}
Earlier drafts carried multiple near-equivalent predicates and exclusion formulations in parallel.
Those variants are useful for technical experimentation, but they obscure claim traceability in the main proof narrative.
In this revision, the main text keeps one canonical predicate and one canonical exclusion statement; alternatives are treated as supporting material only.

\paragraph{Current mathematical status.}
The manuscript remains a proposed proof framework under review.
The RH deduction in the main text is conditional on open theorem targets.
No claim of complete proof is made.

\paragraph{Interpretive boundary.}
Heuristic motivation, computational checks, and architectural planning can guide research, but they are not substitutes for theorem-level proofs.
Formal validity depends only on statements proved under explicit hypotheses.


\section{Limitations and Open Problems}
This section records the present limits of the argument and identifies the exact burden of proof that remains.
It is written for mathematical audit: each unresolved bridge is stated explicitly so that reviewers can test the framework at its most vulnerable points.

\subsection*{Scope and evidentiary standard}
The manuscript proposes a proof framework.
Its intuitive geometric language is intended to guide construction and suggest mechanisms, but it does not by itself constitute proof.
Only statements supported by complete definitions, hypotheses, and logically valid derivations should be treated as formal results.

\subsection*{Primary open bridges}
\begin{enumerate}
  \item \textbf{Rigorous operator construction.}
  The operator/spectral construction must be specified on a precise functional-analytic domain, with well-defined action, closure properties, and spectral meaning.
  At present, this construction is programmatic rather than complete.
  Any downstream claim that depends on the spectrum is therefore conditional on this step.

  \item \textbf{Proof of the zero--balance equivalence.}
  Any claimed equivalence between the proposed balance criteria and nontrivial zeros of $\zeta(s)$ must be proved as a theorem in both directions.
  This equivalence cannot be assumed from analogy, numerics, or partial correspondences.
  Until a full equivalence proof is written, the balance criteria should be treated as candidate reformulations.

  \item \textbf{Rigorous uniqueness of the critical line.}
  The argument requires a proof that the relevant mechanism selects $\Re(s)=\tfrac12$ uniquely, not merely that it is compatible with that line.
  Exclusion of off-line alternatives must be established by explicit estimates or structural theorems under stated hypotheses.
  Without such a proof, uniqueness of the critical line remains open within this framework.
\end{enumerate}

\subsection*{Interpretation of intermediate evidence}
Computational checks, model examples, and heuristic consistency tests are useful for stress-testing definitions and detecting contradictions, but they are not substitutes for proof.
They may justify further investigation; they do not discharge formal obligations.

\subsection*{Policy for unresolved steps}
Any unresolved dependency should be treated as open and cited where it is used.
No conclusion stronger than the proved dependencies is warranted.
Accordingly, the manuscript should be read as materials for mathematical audit of a proposed proof framework, with explicit open bridges rather than implicit assumptions.


% ===== Appendices =====
\appendix

\section{Intuition and Motivation}
This appendix collects the intuitive and motivational material that underlies
the formal development. The material here is \emph{not} part of the proof;
it is provided to help readers understand the heuristic picture that motivates
the definitions and draft theorem targets of the main text.

\subsection{The Balance Metaphor}

Consider a balance scale with a pivot at the center. Objects placed on either
side must be equal in weight for the scale to be in equilibrium. The pivot is
the unique point of balance.

The completed zeta function $\xi(s)$ is symmetric under reflection about
$\Re(s) = \tfrac{1}{2}$: it satisfies $\xi(s) = \xi(1-s)$. In the balance
metaphor, $\xi$ is a scale, and the reflection $s \mapsto 1-s$ swaps the
two sides. The critical line $\Re(s) = \tfrac{1}{2}$ is the pivot.

The informal picture motivating this manuscript is that zeros of $\zeta$
are ``balance points'' — states of equilibrium. If the equilibrium condition
is defined correctly, it might only be achievable at the pivot. In the current
manuscript status, this remains heuristic motivation; the formal program is
to define the equilibrium condition precisely and prove the corresponding
draft theorem targets.

\subsection{The Standing Wave Analogy}

On a vibrating string fixed at both endpoints, standing waves occur at specific
resonant frequencies. At a resonant frequency, the wave pattern reinforces itself
and forms a stable pattern. Off-resonance, the wave cancels and no stable pattern exists.

In the proposed spectral-coherence picture, the nontrivial zeros of $\zeta$ are
proposed to be analogous to resonant frequencies. The spectral operator $H$ plays
the role of the string, and the ``resonance condition'' is the balance criterion.
Heuristically, one expects resonance --- balance --- to concentrate on frequencies
corresponding to the critical line; formally, this expectation is represented by
Target~B in the main results as a theorem target under construction.

This analogy is presented here precisely because it is \emph{only} an analogy.
The formal content is in \S4 and \S8.

\subsection{Why the Critical Line Should Be Special}

The reflection $s \mapsto 1-s$ is a symmetry of the problem. The fixed points
of this symmetry are exactly the points $s = \tfrac{1}{2} + it$, which form the
critical line.

In many physical and mathematical systems, the fixed-point locus of a symmetry
has special properties not shared by other points. A simple example: the midpoint
of an interval $[a, b]$ is the unique point invariant under the reflection that
swaps $a$ and $b$.

The framework suggests that the critical line is special in a specific sense:
it may be the only locus where the balance criterion is satisfied. In the current
manuscript vocabulary, this is the statement of Target~B (uniqueness of balance
locus), Theorem~\ref{thm:target-uniqueness}, which is a draft theorem target and
not yet proved in this manuscript version.

\subsection{A Word on Higher-Dimensional Extensions}

Some versions of the framework under exploration involve lifting the complex variable
$s$ to a quaternionic or higher-dimensional setting. The motivation is that certain
symmetry arguments may become cleaner or more natural in a higher-dimensional space.

These extensions are exploratory and are not currently part of the formal argument.
They are noted here as a direction for future investigation.


\section{Operator Setup and Technical Background}
This appendix collects technical background relevant to the proposed spectral
construction. It is intended as a reference for readers less familiar with
operator theory and spectral analysis.

\subsection{Basic Operator Theory}

Let $\hs$ be a complex Hilbert space with inner product $\langle \cdot, \cdot \rangle$.

\begin{itemize}
  \item A \emph{linear operator} $T \colon \mathcal{D}(T) \to \hs$ is a linear
        map defined on a dense subspace $\mathcal{D}(T) \subset \hs$.
  \item $T$ is \emph{self-adjoint} if $\mathcal{D}(T) = \mathcal{D}(T^*)$ and
        $T = T^*$.
  \item The \emph{spectrum} of $T$, denoted $\sigma(T)$, is the set of $\lambda \in \C$
        for which $T - \lambda I$ is not invertible (as a bounded operator from
        $\hs$ to $\hs$).
  \item A self-adjoint operator has real spectrum: $\sigma(T) \subset \R$.
\end{itemize}

\begin{remark}
  Self-adjointness is a strictly stronger condition than symmetry ($T \subset T^*$).
  Many operator constructions in spectral approaches to RH fail at the step of
  establishing self-adjointness.
\end{remark}

\subsection{Spectral Theorem}

For a self-adjoint operator $T$ on $\hs$, the spectral theorem guarantees the
existence of a projection-valued measure $E$ such that
\[
  T = \int_\R \lambda \, dE(\lambda).
\]
Trace-class operators have a well-defined trace $\mathrm{tr}(T) = \int_\R \lambda \, d\mu_T(\lambda)$
where $\mu_T$ is the spectral measure.

\subsection{Trace Formulas}

The Selberg trace formula and related trace formulas express geometric or spectral
data as a sum over eigenvalues. In approaches to RH via spectral theory, the goal
is to construct an operator whose trace formula reproduces the explicit formula
for $\zeta$ or its logarithmic derivative.

The explicit formula for $\zeta$ relates sums over primes to sums over zeros:
\[
  \sum_p \log p \cdot p^{-s} = -\frac{\zeta'(s)}{\zeta(s)} - (\text{explicit terms}).
\]
A spectral construction matching this would give the zeros as spectral data.

\todo{The relationship between the proposed operator $H$ and any explicit formula
should be worked out explicitly when $H$ is fully defined.}

\subsection{Functional Analysis Prerequisites}

The following concepts may be needed in the construction of $H$:

\begin{itemize}
  \item Hardy spaces $H^2(\C^+)$ on the upper half-plane.
  \item Paley--Wiener theory relating Fourier transform support to analytic continuation.
  \item The theory of Krein--Feller operators or Schrödinger operators on intervals.
  \item Mellin transforms and their relationship to $\zeta$.
\end{itemize}

\todo{As the construction progresses, the specific functional-analytic tools needed
will be identified and their prerequisites will be stated here.}


\section{Computational Notes}
This appendix describes the computational support for the manuscript. The code
and notebooks in \texttt{code/} are provided to visualize and explore the
proposed spectral-coherence structure numerically. Computational results are
illustrative and do not substitute for formal proofs.

\subsection{What the Code Does}

\begin{itemize}
  \item \texttt{code/notebooks/spectral-visualization.ipynb}: Plots the magnitude
        and phase of $\zeta(s)$ along and near the critical line. Visualizes the
        reflection symmetry $s \mapsto 1-s$.
  \item \texttt{code/notebooks/critical-line-geometry.ipynb}: Plots the critical
        line and the placement of known nontrivial zeros. Illustrates the geometric
        intuition of the balance framework.
  \item \texttt{code/src/placeholder.py}: Placeholder module with utility functions
        for computing $\zeta(s)$ numerically (using the \texttt{mpmath} library
        for high-precision arithmetic).
\end{itemize}

\subsection{Numerical Limitations}

All numerical computations of $\zeta(s)$ are approximate. The known zeros have
been computed to many decimal places by other researchers (see \cite{Odlyzko1987,
PlattTrudgian2021}), and the code in this repository does not aim to extend the
numerical verification of RH. Rather, it supports visual understanding of the
framework.

No claim in the manuscript relies on the output of these computations.

\subsection{Reproducing the Computations}

To reproduce the notebooks:

\begin{enumerate}
  \item Install Python 3.9 or later with the packages listed in
        \texttt{code/README.md}.
  \item Navigate to \texttt{code/notebooks/} and launch Jupyter:
        \texttt{jupyter notebook}.
  \item Run all cells in order.
\end{enumerate}

See \texttt{code/README.md} for full instructions.


% ===== References =====
\printbibliography[heading=bibintoc, title={References}]

\end{document}
